\chapter{Multiple Time Scales}

\section{The Fountain Reconsidered}

Chapter 16 left a loose thread on purpose. The fountain was offered as a dynamic fixed point, a region under continuous proposal that nonetheless persists, and the example did real work establishing that stability need not mean inactivity. But the example also quietly assumed something Chapter 16 never examined: that the fountain's water and the fountain's form update on the same clock, so that a single fixed-point equation, checked once per update, could meaningfully describe both at once. This assumption does not survive close inspection. The water occupying the fountain at this instant is entirely different water a minute from now, flowing, splashing, replaced continuously; nothing about its appearance or its flow field is remotely fixed from one moment to the next. The fountain's stone basin, by contrast, is not meaningfully different a minute from now, or an hour from now, or most days. Treating these two facts as though they were answered by the same equation, checked at the same rate, was a simplification, and this chapter exists to undo it.

\section{Fast and Slow Components}

The field's five components do not all change at comparable rates, and this book's supporting technical material already anticipated the division this section makes explicit. Appearance, z, and flow, v, are fast: they are exactly the components most sensitive to whatever is happening right now, water moving through the fountain, light shifting across the wall, and they are proposed against, projected, and committed on a clock fine enough to capture that immediacy. Coherence, Φ, is slow: it aggregates over many fast-timescale events before it meaningfully shifts, since a single moment's flow or appearance rarely changes how stable a region is overall, and the topological or identity-level structure that Φ helps define moves slower still, changing only when enough has accumulated to actually warrant it. Residue, R, sits between the two, accumulating at whatever rate events actually occur, sometimes fast, when a region sees frequent activity, sometimes slow, when little happens for long stretches, and this book will treat R's rate as inherited from whatever it is accumulating rather than fixed to either extreme.

This is not a new mechanism bolted onto the write cycle Chapters 12 through 15 already established. It is the observation that the same cycle, propose, project, commit, does not have to run at one uniform rate across every component it touches. Fast components run through many cycles in the time slow components run through one, and a region can be, quite legitimately, in the middle of vigorous fast-timescale activity while its slow-timescale structure sits still.

\section{Separate Clocks, Same Cycle}

It is worth being precise about what does and does not change once fast and slow components are given separate clocks. Nothing about proposal, projection, or commit is different in kind; a fast-timescale update still proposes a \(\Delta\psi\), still projects it onto the nearest state admissible under C, still either commits or refuses the result exactly as Chapters 13 through 15 described. What changes is only the rate at which this cycle runs for a given component. The fountain's water proposes, projects, and commits new appearance and flow values on a clock fine enough to track individual moments of motion. The fountain's basin proposes, projects, and commits new coherence values on a clock coarse enough that, across any span of time a casual observer would notice, nothing about it visibly changes at all, even though the cycle is, in principle, still running underneath.

\section{Revisiting Fixed Points}

This sharpens the trivial-versus-dynamic distinction Chapter 16 drew, and it is worth returning to that chapter's language directly rather than leaving the two accounts standing side by side unreconciled. A dynamic fixed point, as Chapter 16 described it, is a region under continuous proposal that nonetheless returns to the same state. It is now possible to say more precisely what this means: a dynamic fixed point is fixed at the slow timescale while remaining active, often highly active, at the fast timescale. The fountain's form, Φ, is a genuine fixed point, checked and reconfirmed on the slow clock; the fountain's water is not fixed at all, cycling continuously on the fast clock, and the fountain as a whole is correctly called stable only because these two facts are being reported at two different rates rather than one. A trivial fixed point, by contrast, is fixed at both timescales simultaneously: nothing proposed, nothing cycling, at either rate. The bare, undisturbed paving from Chapter 16 is trivial in this fuller sense, not merely because no single \(\Delta\psi\) was proposed for it at some particular moment, but because neither its fast nor its slow components are undergoing any cycle at all.

\section{When Timescales Interact}

The two clocks are not fully independent, and the most important claim this chapter has to make is that fast-timescale activity, given enough time, is what eventually forces slow-timescale change. A single moment of water striking the fountain's basin does nothing to its coherence; the basin's Φ is far too slow to respond to any individual fast-timescale event. But water striking the same spot, moment after moment, across enough fast-timescale cycles, accumulates residue, R, at that point, exactly as Chapter 11 described, and residue accumulated far enough eventually licenses a slow-timescale proposal that would not have been admissible earlier: erosion, a hairline crack, eventually a change to the basin's coherence itself. This is how a region moves from being a dynamic fixed point to no longer being one, not through any single dramatic event, but through fast-timescale activity slowly accumulating, via residue, until it crosses a threshold the slow timescale is finally forced to answer for. The two clocks are coupled through exactly the mechanism this book spent Chapter 11 establishing, and this is arguably the clearest illustration yet of why residue needed its own component rather than being treated as a byproduct of appearance.

\section{Closing Part III}

This chapter completes the account Part III set out to give. Chapter 12 established that writing is addition rather than replacement, \(M_{t+1}\) = \(M_t\) + \(\Delta\psi\). Chapter 13 gave \(\Delta\psi\) its origin, a combination of the field's own intrinsic tendencies and a learned, agent-driven forcing term, selected among competing candidates by relevance. Chapter 14 gave every proposal somewhere to be checked, projected onto the nearest state the accumulated constraints of C actually permit. Chapter 15 separated that computation from the historical fact of commitment, and gave refusal a proper home in the gap between them. Chapter 16 showed that persistence itself, the very property this book set out from in Chapter 6 to explain, is nothing more exotic than a fixed point of this cycle, and this chapter has shown that the cycle does not run at one rate, but at as many rates as the field's components require, coupled to one another through the same residue that Part II had already given a home.

Nothing in these six chapters, deliberately, has said what C actually is in full, what argmin means when the space being minimized over has the structure this book's fields require, or what guarantees, if any, ensure that a fixed point exists at all rather than being an accident of a particular courtyard's particular history. Those are Part IV's questions. Chapter 18 begins there, with the constraint manifold this Part has invoked six times without yet defining.
